% AUTO-GENERATED by scripts/aggregate_results.py -- do not edit.
\begin{table}[!t]
\centering
\caption{Probing the nesting hierarchy (Eqs.~\eqref{eq:probe}--\eqref{eq:spearman}). Out-of-sample $R^2$ for eleven physiological descriptors as a function of prefix length. Every descriptor improves monotonically to $d$=512, so the saturation criterion of Eq.~\eqref{eq:sat} lands at 256 or 512 for all eleven and has little rank variance to correlate against the hypothesised ordering. The consequential observation is the contrast with Table~\ref{tab:main}: diagnostic AUC is flat from $d$=16 while descriptor recoverability rises by a mean factor of \R{ProbeGain}.}
\label{tab:probe}
\scriptsize
\setlength{\tabcolsep}{3.2pt}
\begin{tabular}{@{}lrrrrrrrrr@{}}
\toprule
& \multicolumn{6}{c}{$R^2_j(m)$} & & & \\
\cmidrule(lr){2-7}
Descriptor & 16 & 32 & 64 & 128 & 256 & 512 & $\times$ & $s_j$ & $c_j$ \\
\midrule
heart rate      & .373 & .554 & .657 & .735 & .789 & .826 & 2.2 & 256 & 16 \\
RR mean         & .306 & .456 & .566 & .634 & .679 & .710 & 2.3 & 256 & 16 \\
RR SD           & .068 & .077 & .099 & .129 & .170 & .191 & 2.8 & 512 & 16 \\
RR RMSSD        & .066 & .077 & .096 & .119 & .158 & .179 & 2.7 & 512 & 16 \\
QRS duration    & .102 & .139 & .159 & .201 & .247 & .296 & 2.9 & 512 & 32 \\
QT interval     & .143 & .188 & .227 & .263 & .289 & .330 & 2.3 & 512 & 32 \\
ST level        & .297 & .321 & .368 & .398 & .421 & .440 & 1.5 & 256 & 64 \\
T amplitude     & .550 & .595 & .607 & .646 & .675 & .713 & 1.3 & 512 & 128 \\
R amplitude     & .444 & .582 & .660 & .742 & .783 & .830 & 1.9 & 512 & 256 \\
P amplitude     & .113 & .165 & .201 & .296 & .402 & .502 & 4.4 & 512 & 256 \\
QRS axis        & .135 & .238 & .342 & .461 & .514 & .583 & 4.3 & 512 & 128 \\
\midrule
\multicolumn{10}{@{}l}{Spearman $\rho(c_j,s_j)=0.399$, $p$=0.224 $\Rightarrow$
hierarchy \textbf{not supported}} \\
\multicolumn{10}{@{}l}{Macro AUC over the same range: 0.9264 $\rightarrow$
0.9251 (\textbf{flat})} \\
\bottomrule
\end{tabular}
\end{table}
